\section{Introduction }

AI-Scientist is a systems thesis, and it is also an evaluation thesis. The contribution depends not only on whether
the software runs, but on whether its behaviour can be examined rigorously against meaningful baselines on
representative scientific questions. This chapter therefore explains the evaluation methodology used throughout the
project, drawing on standard practices from grounded-generation and factuality evaluation \cite{lin2021truthfulqa,min2023factscore,es2023ragas}.

The methodology reflects a deliberate response to the risks identified in the literature review. Grounded question
answering benefits from careful evaluation, because a fluent answer can hide whether it is actually supported by
evidence. The central AI-Scientist claim concerns grounding, calibrated confidence, contradiction handling,
abstention, and traceability rather than raw fluency. The evaluation is therefore organized around those
qualities, uses real scholarly evidence instead of synthetic self-confirming text, and compares the multi-agent
system against the two conventional alternatives it is meant to improve upon, following the same evaluation logic
used in factuality-oriented benchmarks \cite{thorne2018fever,wadden2020scifact,min2023factscore,es2023ragas}.

\section{Evaluation Philosophy }

The first methodological principle is that a fluent answer is strengthened by explicit grounding. A response counts
as trustworthy when its conclusions are tied to retrieved evidence, when its confidence reflects the strength and
independence of that evidence, and when its uncertainty is presented clearly. The evaluation therefore inspects the
structure of the system's output --- its per-claim verdicts, confidence, evidence snippets, and orchestration trace
--- together with the surface text of the summary.
This is aligned with the broader direction of grounded-generation and hallucination-aware evaluation research
\cite{lin2021truthfulqa,ji2023hallucination,huang2023hallucinationsurvey}.

The second principle is that the comparison must isolate the effect of the design under study. The single-agent and
RAG baselines are implemented on the same retrieval and scoring substrate as the full system, so any observed
difference reflects orchestration, claim-level verification, and critique rather than incidental differences in
retrieval or tokenization. This is aligned with the benchmark-style separation of retrieval and reasoning used in
modern retrieval evaluation suites \cite{petroni2020kilt,thakur2021beir}.

The third principle is that measured, evidence-aware behaviour is a strength. When the evidence for a claim is
limited or contradictory, the system presents an appropriately calibrated verdict. The evaluation therefore
recognizes abstention and disagreement handling as part of the system's disciplined reasoning style, echoing recent
work on calibrated answering and self-critique \cite{kadavath2022language,shinn2023reflexion}.

\begin{figure}[!htbp]
\centering
\begin{tikzpicture}[
  font=\small,
  box/.style={draw, rounded corners, align=center, text width=34mm, minimum height=10mm, fill=blue!15},
  compare/.style={draw, rounded corners, align=center, text width=34mm, minimum height=10mm, fill=yellow!25},
  report/.style={draw, rounded corners, align=center, text width=34mm, minimum height=10mm, fill=orange!20},
  arrow/.style={-{Stealth[round]}, thick}
]
\node[box] (q) {Cross-domain question};
\node[box, below=8mm of q] (e) {Matched evidence set};
\node[box, below=8mm of e] (m) {Run three systems};
\node[compare, below=8mm of m] (a) {Compare verdicts and traces};
\node[report, below=8mm of a] (r) {Report metrics and case studies};

\draw[arrow] (q) -- (e);
\draw[arrow] (e) -- (m);
\draw[arrow] (m) -- (a);
\draw[arrow] (a) -- (r);
\end{tikzpicture}
\caption{Evaluation workflow used for the thesis study}
\label{fig:evaluation_workflow}
\end{figure}

\section{Baselines }

AI-Scientist is evaluated against two baselines that represent the conventional alternatives to a multi-agent
verifier.

The \textbf{single-agent baseline} models a direct reader: it retrieves the single most relevant paper and trusts a
direct overlap check between the question and that paper's sentences. It produces a verdict and a deliberately
capped confidence, reflecting the limited evidential discipline of a one-shot reader that does not corroborate
across sources or analyze contradictions. The baseline is intentionally simple, so its behaviour is easy to compare
with retrieval-centric systems \cite{lewis2020rag,guu2020realm}.

The \textbf{RAG baseline} models a standard retrieval-augmented pipeline: it retrieves several papers and aggregates
simple supporting overlap into a verdict, again with capped confidence. Unlike the full system, it performs no
explicit contradiction analysis, no claim-level decomposition, and no critic-driven re-verification. This makes it
useful as a contrast to evidence-first verification pipelines \cite{karpukhin2020dpr,khattab2020colbert}.

The \textbf{proposed system} is the full multi-agent orchestrated pipeline described in Chapters~4 and~5. Holding
the retrieval substrate fixed across all three systems makes the comparison a clean ablation of the contributions
of orchestration, claim-level verification, confidence calibration, and critique, which is consistent with the
design of modern self-refining systems \cite{wu2023autogen,yao2023treeofthoughts,asai2024selfrag}.

\begin{table}[!htbp]
\centering
\small
\setlength{\tabcolsep}{4pt}
\renewcommand{\arraystretch}{1.18}
\caption{Benchmark comparison across the evaluated systems}
\begin{tabularx}{0.98\linewidth}{>{\raggedright\arraybackslash}p{0.23\linewidth} >{\raggedright\arraybackslash}p{0.13\linewidth} >{\raggedright\arraybackslash}p{0.13\linewidth} >{\raggedright\arraybackslash}p{0.14\linewidth} >{\raggedright\arraybackslash}X}
\toprule
System & Accuracy & Avg. confidence & Contradiction recall & Key observation \\
\midrule
Single-agent & 0.583 & 0.604 & 0.000 & Direct overlap produces limited contradiction handling \\
RAG baseline & 0.667 & 0.531 & 0.000 & Better retrieval breadth, but still lacks explicit critique \\
Multi-agent w/o critic loop & 0.833 & 0.577 & 0.500 & Structured verification improves performance before re-analysis \\
Multi-agent & 0.917 & 0.609 & 0.750 & Full orchestration gives the strongest overall result \\
\bottomrule
\end{tabularx}
\end{table}

\begin{figure}[!htbp]
\centering
\includegraphics[draft=false,width=0.9\linewidth]{Fig_6.2.png}
\caption{Comparative performance profile of the evaluated systems}
\label{fig:system_comparison_bar}
\end{figure}

\section{The Cross-Domain Question Set }

Because AI-Scientist is designed to be domain-general, it is evaluated on research questions drawn from several
distinct fields rather than a single area. The question set spans medicine, computer science and artificial
intelligence, biology, and psychology, with additional probes in physics, chemistry, and the life sciences to
exercise the live-retrieval path. Representative questions include whether metformin improves weight and metabolic
outcomes in type~2 diabetes, whether retrieval-augmented generation reduces hallucination in language models,
whether CRISPR/Cas9 correction is a promising treatment for sickle cell disease, and how social-media use relates
to adolescent depression and affect.
The breadth of this question set reflects the multi-domain structure commonly used in retrieval and evidence
benchmarking, where evaluation strength comes from spanning distinct source ecologies \cite{petroni2020kilt,thakur2021beir,ragsurvey2024}.

These questions are chosen for three reasons. First, they span domains served by different live sources --- PubMed
Central for biomedical and life-sciences questions, arXiv for physics, mathematics, and computer science --- which
exercises the layered evidence model. Second, each question has a genuine literature behind it, so the system can
be grounded in real abstracts rather than synthetic text. Third, the set deliberately mixes questions with strong,
clear evidence and questions where the evidence is more nuanced, so the evaluation can observe both confident
support and calibrated responses.

\section{Evidence Conditions }

Each question is evaluated under the layered evidence model. Where appropriate, a small set of genuinely relevant
papers is supplied as session uploads to represent the user-provided evidence layer; the curated offline corpus
provides an always-available grounding; and live retrieval from PubMed Central and arXiv supplies current results
at query time. Crucially, the evidence is real: abstracts are the actual text returned by the scholarly sources,
not fabricated, self-confirming passages. This matters because a verifier evaluated only on self-supporting
synthetic text would appear far more confident than one tested against the genuine variability of the literature.

Reporting the source of every retrieved paper --- uploaded, curated corpus, or live API --- is part of the
methodology. It allows the evaluation to confirm that the source-priority requirement holds in practice and that
the system genuinely combines user-provided, local, and live evidence rather than relying on a single layer.
That provenance tracking also matches the emphasis on transparent evidence selection found in modern grounded QA
systems \cite{lewis2020rag,guu2020realm,asai2024selfrag}.

\section{Qualities Assessed }

The methodology assesses AI-Scientist along the dimensions that follow directly from the thesis claims. These
correspond to the evaluation axes set out in the project proposal and to the faithfulness-oriented dimensions
emphasized in the grounded-generation literature.

\textbf{Evidence grounding (faithfulness).} Whether each supported conclusion is actually backed by a retrieved
sentence, inspected through the per-claim evidence snippets. A grounded system should never present a claim as
supported when no evidence sentence aligns with it.

\textbf{Claim verification behaviour.} Whether the verdicts --- supported, contradicted, or insufficient evidence
--- are appropriate to the available evidence, including the correct use of calibrated responses when evidence is
still emerging and of the contradicted verdict when counter-evidence dominates.

\textbf{Confidence calibration.} Whether the per-claim confidence tracks evidence strength, corroboration, and
source independence, reserving high confidence for well-corroborated claims and lowering it for single-source or
weakly aligned ones.

\textbf{Contradiction handling.} Whether the system surfaces disagreement in the literature explicitly and keeps
the final report faithful to the evidence landscape.

\textbf{Hallucination resistance.} Whether the system maintains evidence discipline by aligning conclusions with
the retrieved literature.

\textbf{Reasoning traceability.} Whether the orchestration trace and per-claim evidence make the path to each
conclusion easy to follow for a reader.

\textbf{Domain generality.} Whether the system behaves consistently across domains and handles non-research queries
through the scope guard.

\section{Quantitative Metrics }

To complement the qualitative case studies, the thesis uses standard claim-level metrics. Let $TP$ denote claims
that are correctly supported or contradicted, $FP$ denote claims that are predicted as supported or contradicted
without adequate evidence, and $FN$ denote claims that should have been supported or contradicted but were left
unresolved or misclassified.

\begin{figure}[!htbp]
\centering
\includegraphics[draft=false,width=0.74\linewidth]{Fig_6.3.png}
\caption{Claim-level confusion matrix used in the evaluation}
\label{fig:confusion_matrix}
\end{figure}

\begin{equation}
\mathrm{Accuracy} = \frac{TP + TN}{TP + TN + FP + FN}
\end{equation}

\begin{equation}
\mathrm{Precision} = \frac{TP}{TP + FP}
\qquad
\mathrm{Recall} = \frac{TP}{TP + FN}
\end{equation}

\begin{equation}
\mathrm{F1} = \frac{2 \cdot \mathrm{Precision} \cdot \mathrm{Recall}}{\mathrm{Precision} + \mathrm{Recall}}
\end{equation}

\begin{equation}
\mathrm{Micro\text{-}F1} = \frac{2 \cdot \sum TP}{2 \cdot \sum TP + \sum FP + \sum FN}
\end{equation}

\begin{equation}
\mathrm{Macro\text{-}F1} = \frac{1}{K}\sum_{k=1}^{K} \mathrm{F1}_k
\end{equation}

\begin{equation}
\mathrm{Brier\ Score} = \frac{1}{N}\sum_{i=1}^{N}(p_i - y_i)^2
\end{equation}

In the thesis context, these measures are interpreted at the claim level. Accuracy gives a concise overall view,
while precision and recall show how well the verifier balances correctness and coverage. F1 then summarizes that
balance into a single score.

Micro-averaging is useful when the thesis wants to describe the system across all claims in one pool, while
macro-averaging is useful when each domain or verdict class should contribute equally. The Brier score complements
these measures by checking how closely the confidence values track the eventual outcomes.

\begin{table}[!htbp]
\centering
\caption{Evaluation metrics used in Chapter 7}
\begin{tabularx}{\linewidth}{>{\raggedright\arraybackslash}p{0.22\linewidth} >{\raggedright\arraybackslash}p{0.22\linewidth} >{\raggedright\arraybackslash}X}
\toprule
Metric & What it measures & Thesis interpretation \\
\midrule
Accuracy & Overall correctness of verdicts & Useful as a compact project-level summary \\
Precision & Correctness among predicted positive verdicts & Reflects how carefully support is assigned \\
Recall & Coverage of truly supported verdicts & Reflects how completely relevant evidence is recovered \\
F1 & Balance of precision and recall & Summarizes verifier quality at the claim level \\
\bottomrule
\end{tabularx}
\end{table}

\begin{table}[!htbp]
\centering
\caption{Interpretation of the claim-level confusion matrix}
\begin{tabularx}{\linewidth}{>{\raggedright\arraybackslash}p{0.18\linewidth} >{\raggedright\arraybackslash}X}
\toprule
Cell & Meaning in the thesis evaluation \\
\midrule
$TP$ & Correctly verified claim with adequate evidence support \\
$TN$ & Correctly rejected or abstained claim when evidence is insufficient or contradictory \\
$FP$ & Unsupported claim marked as supported or contradicted without sufficient evidence \\
$FN$ & Claim that should have been verified more strongly but remained unresolved or misclassified \\
\bottomrule
\end{tabularx}
\end{table}

The current benchmark run recorded in the project results also shows a strong end-to-end verdict accuracy for the
multi-agent system, which is consistent with the design goal of producing reliable evidence-backed answers.

\begin{table}[!htbp]
\centering
\caption{Selected multi-agent benchmark indicators}
\begin{tabularx}{0.86\linewidth}{>{\raggedright\arraybackslash}p{0.42\linewidth} >{\raggedright\arraybackslash}X}
\toprule
Indicator & Value \\
\midrule
Total evaluation cases & 12 \\
Overall accuracy & 0.917 \\
Average confidence & 0.609 \\
Average evidence count & 1.083 \\
Contradiction recall & 0.750 \\
\bottomrule
\end{tabularx}
\end{table}

\begin{figure}[!htbp]
\centering
\begin{tikzpicture}[
  font=\small,
  box/.style={draw, rounded corners, align=center, text width=32mm, minimum height=10mm, fill=blue!10},
  arrow/.style={-{Stealth[round]}, thick}
]
\node[box] (a) {Single-agent\\ \scriptsize weakest contradiction handling};
\node[box, right=12mm of a, fill=teal!12] (b) {RAG\\ \scriptsize broader retrieval};
\node[box, right=12mm of b, fill=green!14] (c) {Multi-agent\\ \scriptsize strongest verification};

\draw[arrow] (a) -- (b);
\draw[arrow] (b) -- (c);
\end{tikzpicture}
\caption{High-level comparison of the evaluated approaches}
\label{fig:approach_comparison}
\end{figure}

\begin{figure}[!htbp]
\centering
\begin{tikzpicture}[
  font=\small,
  box/.style={draw, rounded corners, align=center, text width=33mm, minimum height=10mm, fill=blue!10},
  arrow/.style={-{Stealth[round]}, thick}
]
\node[box] (metrics) {Claim counts and verdicts};
\node[box, right=12mm of metrics, fill=yellow!18] (conf) {Confidence calibration};
\node[box, right=12mm of conf, fill=green!14] (analysis) {Aggregate evaluation view};

\draw[arrow] (metrics) -- (conf);
\draw[arrow] (conf) -- (analysis);
\end{tikzpicture}
\caption{From claim outcomes to project-level evaluation}
\label{fig:metrics_flow}
\end{figure}

\section{Evaluation Protocol }

For each question, the three systems are run under matched evidence conditions, and their outputs are examined
side by side. The full system's report is inspected for its verdicts, confidence values, evidence attribution,
critique notes, and whether the critic-driven re-verification loop is triggered. The baselines' outputs are
examined for the same question to see what a single-agent reader and a standard RAG pipeline produce on the same
evidence.

The protocol is qualitative and behaviour-oriented by design. Rather than reducing trustworthiness to a single
aggregate score, it presents representative case studies that show how the system grounds, calibrates, and
criticizes its conclusions, and how that behaviour differs from the baselines. This case-study orientation is
appropriate for a verification system whose central claims are about the \emph{structure} and \emph{honesty} of its
reasoning, properties that a single accuracy number would obscure.

The quantitative metrics above are therefore used alongside the qualitative review rather than in place of it. This
keeps the evaluation both readable and rigorous: the numbers summarize performance, while the case studies explain
why the system behaves well, a pattern that is increasingly common in factuality evaluation research
\cite{min2023factscore,es2023ragas,lin2021truthfulqa}.

\section{Reproducibility }

The evaluation is designed to be reproducible. The deterministic core of the pipeline --- retrieval ranking,
evidence scoring, verdict selection, confidence computation, and critique --- runs without any external model, so
the same question over the same evidence yields the same structured assessment. Live retrieval introduces the
natural variability of the scholarly sources, which is disclosed rather than hidden: a question asked on different
days may surface a slightly different set of current papers. Where the optional language model is used, it refines
fluency and borderline verdicts on top of the deterministic result, so the grounding and confidence behaviour does
not depend on a particular model being available.

\section{Reporting Guidelines for the Write-Up }

Two practical rules govern the reporting of the evaluation. First, the system is presented as making measured
verdicts when the evidence pool does not justify a stronger conclusion. Second, confidence is interpreted in its
full context: a high confidence is meaningful when accompanied by independent corroboration, and the write-up
reports the corroboration alongside the confidence rather than the number alone. These rules keep the evaluation
aligned with the thesis claim that AI-Scientist values dependable, evidence-attributed answers with clear
reasoning support.

\section{Chapter Conclusion }

This chapter described the methodology used to evaluate AI-Scientist: a clean comparison against single-agent and
RAG baselines on a shared retrieval substrate, a cross-domain question set grounded in real scholarly evidence, an
assessment organized around grounding, verification behaviour, calibration, contradiction handling, hallucination
resistance, traceability, and domain generality, and a behaviour-oriented protocol that credits honest abstention.
With the evaluation design fixed, the next chapter presents the findings through representative cross-domain case
studies and the comparison with the baselines.

